\begin{textsample}{POS Dim 7 – global\_south\_2024\_11 – Score 3.00 – global\_south\_2024\_11/117295995.txt}  \label{ex:f7_pos_012}
Why Is Malaysia Seeking Israel ' s Expulsion From The United Nations?
Malaysia is planning to seek the backing of Arab nations at an emergency summit next week to support a proposal for Israel ' s expulsion from the United Nations, as Prime Minister Anwar Ibrahim informed lawmakers.
Malaysia is planning to seek the backing of Arab nations at an emergency summit next week to support a proposal for Israel ' s expulsion from the United Nations, as Prime Minister Anwar Ibrahim informed lawmakers.
He claimed that Israel ' s continued actions in Gaza contravene international law and United Nations directives.
Israel at loggerheads with UNRWA The conflict, which intensified after a Hamas attack in October last year that killed 1,200 Israelis, has since led to the deaths of over 43,000 Palestinians, mainly women and children, and left much of Gaza in ruins.
In a recent development, Israel formally ended its long-standing cooperation agreement with the United Nations Relief and Works Agency ( UNRWA ), accusing it of being infiltrated by Hamas.
This decision @ @ @ @ @ @ @ @ @ @ of humanitarian aid for Palestinians.
Malaysia, outspoken critic of Israel Malaysia, a majority-Muslim nation, has been an outspoken critic of Israel ' s actions in Gaza, including banning the Israeli shipping company Zim, while Malaysians have actively boycotted businesses seen as connected to the United States, Israel ' s primary ally.
Anwar, who has been leading the opposition, reiterated Malaysia ' s commitment to supporting the Palestinian cause for an independent state, condemning Israel ' s actions as acts of"dispossession, terrorism, cruelty, and genocide."He mentioned that Malaysia is considering a proposal to expel Israel from the UN, citing its disregard for laws, principles, and UN decisions.
Malaysia to present proposal at Arab League ' s emergency summit Anwar plans to present this proposal at the Arab League ' s emergency summit in Riyadh next Monday, which \textbf{Saudi} Arabia confirmed will focus on Israel ' s"crimes and violations"in Gaza and its escalated strikes in Lebanon.
According to Huwaida Arraf, a Palestinian-American rights activist @ @ @ @ @ @ @ @ @ @ UN requires evidence of persistent violations of international law and UN decisions, as per articles 5 and 6 of the UN Charter.
However, Arraf pointed out that Security Council approval is essential, which may be challenging due to the United States ' \textbf{veto} power, a right it has frequently exercised to block UN \textbf{resolutions} critical of Israel.
Call for unity among Arab and Muslim countries Anwar underscored the necessity for unity among Arab and Muslim countries in urging the UN to impose a ceasefire and facilitate the delivery of medical and humanitarian aid to Palestinians.
He further suggested establishing a UN Special Committee Against Apartheid to address Israel ' s actions.
Malaysia will also strive for a collective response against Israel ' s recent parliamentary decision to prohibit UNRWA activities within its controlled areas, while pushing for the agency ' s full access to deliver aid to Palestinians.
Israel, on Tuesday, terminated a 1967 agreement that allowed UNRWA ' s operations with diplomatic immunity within Israel, the West Bank, and Gaza.
Additionally, Israeli legislation was @ @ @ @ @ @ @ @ @ @ the agency as a"terrorist"organization.

% matched lemmas: resolution, saudi, veto
\end{textsample}
